\begin{textsample}{POS Dim 4 – 2022\_c5 – Score -1.00 – 2022\_c5/enade\_2022\_tec\_gest\_com\_7f0e07f670.txt}  \label{ex:f4_pos_287}
CAPÍTULO 6 DISTRIBUIÇÃO DE CONCEITOS A finalidade deste capítulo é apresentar um panorama nacional da distribuição do Conceito Enade 2022.
Para traçar esse panorama, recorre -se a tabelas e gráficos, que são acompanhados de comentários e análises que contribuem para a contextualização dos dados neles expostos.
Para visualizar o panorama do Enade/2002 em nível nacional, na primeira seção do capítulo, os dados sobre o número de cursos participantes são organizados por grande região.
Na segunda seção, o mapeamento do conceito do exame é feito por categoria administrativa das IES e modalidade de oferta dos cursos, de maneira cruzada às grandes regiões.
Na última seção, a lógica de apresentação dos dados é feita por organização acadêmica e, também, grande região de funcionamento dos cursos, no país como um todo. 6.1 PANORAMA NACIONAL DA DISTRIBUIÇÃO DOS CONCEITOS A lógica adotada na apresentação dos dados, nesta seção e nas duas seguintes, pressupõe o registro de duas informações de natureza reiterativa.
A primeira informação diz respeito aos valores possíveis do Enade, que variam em níveis de 1 a 5, sendo que, à medida que esses valores aumentam, o curso recebe uma avaliação melhor ; além disso, os cursos com somente um ou sem nenhum concluinte participante não obtêm o Conceito Enade, ficando Sem Conceito ( SC ), como já ressaltado neste relatório.
O segundo registro é que a diferença entre os cursos tabulados neste capítulo e no Capítulo 3 corresponde aos cursos que não foram avaliados, ou seja, àqueles, em princípio, sem estudantes concluintes que tivessem sido inscritos para a prova.
Na Tabela 6.1, são apresentadas a quantidade e a distribuição de cursos de Tecnologia em Gestão Comercial participantes do Enade/2022, por faixa de conceito e grande região.
Observando -se os dados da Tabela 6.1, nota -se que, dos 265 cursos participantes, 96 ( 36,2 \% ) classificaram -se com conceito 3, o valor modal.
Esse foi também o conceito modal em todas as regiões, exceto na região Norte.
O conceito 4 foi o segundo mais frequente em nível nacional ( 21,5 \%, correspondendo a 57 cursos ). 37 cursos ( 14 \% ) classificaram -se com o conceito 2.
Houve, ainda, 10 cursos ( 3,8 \% ) que receberam o conceito 5, e 3 cursos ( 1,1 \% ) que receberam o conceito 1.
Dos 265 cursos de Tecnologia em Gestão Comercial, 62 ( 23,4 \% ) ficaram Sem Conceito ( SC ).
A região Nordeste participou com 49 cursos.
Nessa região, 22 cursos ( 44,9 \% em termos regionais ) obtiveram conceito 3.
O conceito 1 não foi atribuído a nenhum curso e o conceito 2 foi atribuído a 3 cursos ( 6,1 \% ).
Ainda nessa região, 9 cursos ( 18,4 \% ) foram avaliados com o conceito 4, e um curso ( 2 \% ) foi avaliado com o conceito 5. 14 cursos ( 28,6 \% ) ficaram Sem Conceito ( SC ).
Para os 109 cursos participantes da região Sudeste, o conceito 3 foi o conceito modal, como já anteriormente assinalado, tendo sido obtido por 33 cursos ( 30,6 \% ).
O conceito 1 foi atribuído a apenas 2 cursos ( 1,8 \% ), e o conceito 2, a 17 cursos ( 15,6 \% ). 29 cursos ( 26,6 \% ) receberam o conceito 4 e cinco cursos receberam o conceito 5 ( 4,6 \% ).
Nessa região, 23 cursos ficaram Sem Conceito ( SC ).
Em relação aos 72 cursos da região Sul, houve predominância do conceito 3 ( conceito modal ), atribuído a 40,3 \% dos cursos da região, ou seja, atribuído a 29 dos 72 cursos participantes na região Sul.
O conceito 1 não foi atribuído a nenhum curso ; o conceito 2, a 8 cursos ( 11,1 \% ) ; e os conceitos 4 e 5, a 15 ( 20,8 \% ) e a 4 cursos ( 5,6 \% ), respectivamente.
Na região Sul, 16 cursos ficaram Sem Conceito ( SC ).
Já dos 22 cursos participantes na região Centro-Oeste, 10 ( 45,5 \% em termos regionais ) receberam conceito 3.
Um curso ( 4,5 \% ) obteve o conceito 1, 6 cursos ( 27,6 \% ) obtiveram o conceito 2, um curso ( 4,5 \% ), o conceito 4, e nenhum curso alcançou o conceito 5.
Nessa região, 4 cursos ficaram Sem Conceito ( SC ).
Em um gráfico de distribuição cumulativa, em cada valor de abscissa ( eixo x ), o valor apresentado no eixo das ordenadas ( eixo y ) é o valor acumulado da variável ( Conceito Enade ).
Em linhas gerais, em um gráfico desta mesma tipologia, poligonais mais à esquerda representam, no caso dos dados apresentados nesta seção, as grandes regiões com a pior distribuição, e poligonais mais à direita, as grandes regiões com os melhores conceitos.
A partir destas diretrizes, no Gráfico 6.1, é apresentada a distribuição cumulativa do Conceito Enade, segundo a grande região.
No Gráfico 6.1, poligonais mais à esquerda representam valores mais baixos do Conceito Enade, e poligonais mais à direita, valores mais altos. 6.2 CONCEITOS POR CATEGORIA ADMINISTRATIVA, MODALIDADE DE OFERTA DOS CURSOS E GRANDE REGIÃO Na Tabela 6.2, é apresentada a distribuição dos cursos de Tecnologia em Gestão Comercial participantes do Enade/2022, por categoria administrativa e modalidade de oferta, de acordo com os conceitos por eles alcançados, segundo as grandes regiões brasileiras.
De acordo com as informações da Tabela 6.2, dos 265 cursos participantes, 23 eram ministrados em instituições públicas, e 242, em instituições privadas.
Quanto à modalidade de oferta, 140 cursos eram a distância, e 125 eram presenciais.
Ainda conforme as informações da mesma tabela, em termos nacionais, nas instituições públicas, observa -se que nenhum curso foi avaliado com conceito 5.
Dos 23 cursos participantes de IES públicas, o conceito 4 foi o valor modal, atribuído a 11 cursos.
Dos demais cursos de instituições públicas participantes, um curso foi avaliado com o conceito 1.
Nessa categoria, nenhum curso ficou Sem Conceito ( SC ).
Na rede privada, o conceito modal foi o 3, atribuído a 88 cursos.
Dos demais cursos participantes, 2 receberam o conceito 1, 34 receberam o conceito 2, 46 receberam o conceito 4 e 10 cursos receberam o conceito 5.
Nessa categoria administrativa, 62 cursos ficaram Sem Conceito ( SC ).
Dos 125 cursos participantes da modalidade presencial, o conceito 3 foi atribuído a 34 cursos.
Dos demais cursos nessa modalidade de oferta, 2 cursos receberam o conceito 1 ; 17 cursos receberam o conceito 2 ; 30 cursos receberam o conceito 4 e 6 cursos, o conceito 5.
Nessa modalidade de oferta, 16 cursos ficaram Sem Conceito ( SC ).
Dos 140 cursos participantes da modalidade a distância, o conceito 3 foi o valor modal, atribuído a 62 cursos.
Dos demais cursos nessa modalidade de oferta, 1 recebeu o conceito 1 ; 20 cursos receberam o conceito 2 ; 27 cursos, o conceito 4 e 4 cursos, o conceito 5.
Nessa modalidade de oferta, 26 cursos ficaram Sem Conceito ( SC ).
Na região Nordeste, a rede privada concentrou 45 dos 242 cursos participantes.
Entre os cursos oferecidos por IES privadas no Nordeste, 1 curso foi avaliado com conceito 2.
Nenhum curso obteve o conceito 1 ; 20 cursos obtiveram o conceito 3 ; 9, o conceito 4 e 1 curso obteve o conceito 5.
Nessa categoria administrativa e grande região, 14 cursos ficaram Sem Conceito ( SC ).
As instituições públicas da região Nordeste participaram com 4 cursos, 2 com o conceito 2 e 2 com o conceito 3.
Nenhum curso ficou Sem Conceito ( SC ) nessa categoria.
Dos cursos oferecidos na região Nordeste, 23 eram presenciais : 11 ficaram Sem Conceito ( SC ) e 7 ficaram com o conceito 3, o conceito modal.
Na modalidade a distância, eram 23 cursos, tendo 1 obtido o conceito 2 ; 15, o conceito 3 ; 4, o conceito 4 e nenhum, o conceito 5.
Nessa região, 3 cursos ficaram Sem Conceito ( SC ).
Na região Sudeste, a proporção de cursos da rede privada correspondeu a 96 dos 242 cursos participantes.
Nessa combinação de categoria administrativa e grande região, o conceito modal foi o 3, atribuído a 29 cursos. 23 cursos ficaram Sem Conceito ( SC ).
Os demais foram avaliados com os seguintes conceitos : 1 ( 2 cursos ), 2 ( 17 cursos ), 4 ( 20 cursos ) e 5 ( 5 cursos ).
Dos 13 cursos oferecidos por instituições públicas na região Sudeste, o conceito modal foi o 4, atribuído a 9 cursos.
Dos demais cursos, 4 receberam o conceito 3.
Nessa combinação de categoria administrativa e grande região, nenhum curso ficou Sem Conceito ( SC ), e nenhum recebeu o conceito 1, 2 ou 5.
Também na região Sudeste, grande parcela dos cursos foi oferecida na modalidade presencial : 51 de 125 cursos.
Para esses, o conceito 2 foi atribuído a 7 cursos ; o conceito 3 foi atribuído a 11 cursos e o conceito 4 foi atribuído a 17 cursos.
Ficaram Sem Conceito ( SC ) 13 cursos presenciais na região Sudeste.
Na região Sudeste, 58 cursos foram oferecidos na modalidade a distância, tendo sido um curso com conceito 1 ; 10 com conceito 2 ; 22 com conceito 3, o valor modal ; 12 com o conceito 4 e 3 no conceito 5.
Na região Sudeste, 10 cursos a distância ficaram Sem Conceito ( SC ).
As instituições privadas concentraram 70 dos 72 cursos participantes da região Sul.
Desses, 28 cursos obtiveram conceito 3, o conceito modal.
Nessa combinação de categoria administrativa e grande região, nenhum curso recebeu o conceito 1 ; 8 receberam o conceito 2 ; outros 14, o conceito 4 ; e ainda 4 cursos receberam o conceito 5. 16 cursos ficaram Sem Conceito ( SC ).
As instituições públicas na região Sul participaram com apenas 2 cursos, 1 no conceito 4 e 1 no conceito 3.
Nenhum curso da região Sul oferecido por IES públicas ficou Sem Conceito ( SC ).
Na região Sul, 33 dos cursos foram oferecidos na modalidade presencial.
Para esses, a moda consistiu no conceito 3, com 12 cursos.
Os demais cursos presenciais da região foram alocados nos conceitos 2 ( 5 cursos ), 4 ( 6 cursos ), e 5 ( 3 cursos ). 7 cursos ficaram Sem Conceito ( SC ).
Na modalidade a distância, foram oferecidos 39 cursos, nenhum curso com conceito 1 ; 3 com conceito 2 ; 17 com conceito 3, o valor modal ; 9 com o conceito 4 e 1 com o conceito 5.
Observa -se que 9 cursos ficaram Sem Conceito ( SC ).
Na região Centro-oeste, 21 dos 22 cursos participantes eram de IES privadas.
Desses, 6 cursos receberam o conceito 2, 10 receberam o conceito 3 ; 1 recebeu o conceito 4 e nenhum curso recebeu o conceito 5 ou ficou Sem Conceito ( SC ).
Quanto ao único curso participante, oferecido por instituições públicas na região Centro-oeste, apresentou o conceito 1.
Nenhum curso ficou Sem Conceito ( SC ) nessa região.
Também na região Centro-Oeste, a maioria dos cursos foi a distância : 11.
Esses foram alocados nos conceitos 2 ( 4 cursos ), 3 ( 6 cursos ).
Nenhum curso a distância da região ficou Sem Conceito ( SC ) ou nos conceitos 1, 4 ou 5.
Para o curso presencial, 4 participantes foram atribuídos ao conceito 3, o conceito modal. 3 cursos presenciais da região ficaram Sem Conceito ( SC ).
À título de síntese, no Gráfico 6.2, apresenta -se a distribuição do Conceito Enade, segundo a categoria administrativa da IES.
No Gráfico 6.2, como indicado anteriormente, poligonais mais à esquerda representam valores mais baixos do Conceito Enade, e poligonais mais à direita, valores mais altos.
É importante observar que, para cada poligonal presente nesse gráfico, a porcentagem constante no eixo y se refere ao universo de cursos oferecidos por instituições públicas ( poligonal vermelha ) ou privadas ( poligonal azul ). 6.3 CONCEITOS POR ORGANIZAÇÃO ACADÊMICA E GRANDE REGIÃO Na Tabela 6.3, encontra -se a distribuição dos conceitos atribuídos aos cursos participantes do Enade/2022, na área de Tecnologia em Gestão Comercial, por organização acadêmica, segundo as grandes regiões brasileiras.
Como mostram os dados da Tabela 6.3, dos 265 cursos de Tecnologia em Gestão Comercial participantes, 87 eram oferecidos em universidades, 113 em centros universitários, 57 em faculdades e 8 em CEFET/IF.
De acordo com os dados apresentados na mesma tabela, 5 de todos os 10 cursos avaliados com o conceito 5 eram vinculados a universidades.
Os cursos vinculados a esse tipo de organização acadêmica tiveram os conceitos 2 e 3 atribuídos a 48 cursos.
Os demais cursos ofertados por universidades receberam os conceitos 4 ( 17 cursos ) e o conceito 5 ( 5 cursos ). 16 cursos ficaram Sem Conceito ( SC ).
Dos 113 cursos vinculados a centros universitários, o conceito modal também foi o 3, atribuído a 41 cursos. 27 cursos vinculados a esse tipo de organização acadêmica ficaram Sem Conceito ( SC ), e os outros cursos receberam os conceitos 1 ( 1 curso ), 2 ( 19 cursos ), 4 ( 22 cursos ) e 5 ( 3 cursos ).
Dos 57 cursos mantidos por faculdades, 4 receberam conceito 2 ; 17 receberam o conceito 3 ; e 15 receberam o conceito 4.
Dos cursos oferecidos em faculdades, 19 ficaram Sem Conceito ( SC ).
Dos 8 cursos mantidos por CEFET/IF, 1 recebeu o conceito 1, 1 recebeu o conceito 2, 3 receberam os conceitos 3 e 4.
Nenhum curso ficou Sem Conceito ( SC ) ou recebeu o conceito 5.
Considerando -se, separadamente, as regiões brasileiras, verifica -se que, na região Norte, as universidades participaram com 87 cursos avaliados. 16 Sem Conceito ( SC ), 1 no conceito 1, 13 no conceito 2, 35 no 3, 17 no 4 e 5 no conceito 5.
Na região Nordeste, as universidades participaram com 12 dos 87 cursos da área de Tecnologia em Gestão Comercial oferecidos.
O conceito 2 foi atribuído a 1 curso ; o conceito 3 foi atribuído a 6 cursos e o conceito 5 foi atribuído a 1 curso.
Apenas um curso oferecido por universidades no Nordeste ficou Sem Conceito ( SC ).
Os centros universitários contaram com 22 cursos participantes na região Nordeste.
O conceito 2 foi atribuído a 1 curso ; e o conceito 5 não foi atribuído a nenhum curso.
Seis cursos oferecidos por centros universitários no Nordeste ficaram Sem Conceito ( SC ).
As faculdades foram representadas por 14 cursos na região Nordeste, que receberam os conceitos 2 ( 1 curso ), 3 ( 3 cursos ) e 4 ( 3 cursos ). 7 cursos ficaram Sem Conceito ( SC ) nessa combinação de região e organização acadêmica.
Os CEFET/IF foram representados por 1 curso na região Nordeste, que recebeu o conceito 3.
Nenhum curso ficou Sem Conceito ( SC ) nessa combinação de região e organização acadêmica.
Na região Sudeste, as universidades concentraram 39 dos 87 cursos de Tecnologia em Gestão Comercial da região.
Dos cursos oferecidos em IES com este tipo de organização acadêmica, nessa região, o conceito 2 foi atribuído a 8 cursos ; o conceito 3 foi atribuído a 14 cursos e o conceito 4 foi atribuído a 6 cursos.
Sete cursos ficaram Sem Conceito ( SC ).
Os centros universitários participaram com 46 cursos na região Sudeste, dos quais 14 obtiveram o conceito modal 3.
O conceito 1 foi atribuído a um curso ; o conceito 2, a 8 cursos ; o conceito 4, a 11 cursos.
Dois cursos ficaram Sem Conceito ( SC ) nessa combinação de organização acadêmica e região.
Já as faculdades foram representadas por 22 cursos na região Sudeste, dos quais 6 ficaram Sem Conceito ( SC ).
O conceito modal 4 foi atribuído a 10 cursos.
Aos demais cursos foram atribuídos os conceitos 2 ( 1 curso ) e 3 ( 5 cursos ).
Houve 2 cursos participantes de CEFET/IF, na região Sudeste, ambos com conceito 4.
Dos 26 cursos da região Sul, 31 eram mantidos por universidades, para os quais o conceito modal foi 3, atribuído a 12 cursos.
Quatro cursos receberam o conceito 2 e nenhum curso recebeu o conceito 1.
Em IES com esse tipo de organização acadêmica, 7 cursos da região Sul ficaram Sem Conceito ( SC ).
Dos 26 cursos participantes de centros universitários da região Sul, 6 ficaram Sem Conceito ( SC ), e 10 cursos receberam conceito 3, o valor modal.
Os conceitos restantes ficaram assim distribuídos : nenhum curso ( conceito 1 ) ; 4 cursos ( conceito 2 ) ; 7 cursos ( conceito 4 ) ; e um curso ( conceito 5 ).
Houve somente 1 curso vinculado a CEFET/IF na região Sul, que recebeu o conceito 3.
Na região Centro-Oeste, 2 dos 22 cursos oferecidos eram mantidos por universidades.
Para esse tipo de organização acadêmica, os conceitos 1, 2, 4 e 5 não foram atribuídos a nenhum curso.
Dois cursos da região nessa organização acadêmica ficaram com conceito 3.
De maneira semelhante à forma como se procedeu, n em relação à seção anterior, encerra -se esta seção com o Gráfico 6.3, que apresenta a distribuição do Conceito Enade dos cursos de Tecnologia em Gestão Comercial segundo a organização acadêmica da IES.
A leitura desse último gráfico é semelhante à leitura feita dos anteriores, seguindo -se a lógica de que poligonais deslocadas para a esquerda correspondem a valores mais baixos do Conceito Enade ( 1, 2 e 3 ) e poligonais deslocadas para a direita correspondem a valores mais altos do Conceito Enade.

% matched lemmas: 
\end{textsample}
